\chapter{Spherepop: Operator Semantics and Worked Examples}
\label{app:spherepop}

\noindent Spherepop is an irreversible event calculus defined by four primitive
operators acting on a possibility space. This appendix develops the operational
semantics of all four operators, their type signatures, composition rules, and
worked examples demonstrating how Spherepop sequences implement cognitive and
social operations.

\section{The Possibility Space}

\begin{definition}{Possibility Space}
A \emph{possibility space} $\mathcal{P}$ is a partially ordered set of
configurations $\{p_1, p_2, \ldots\}$ with a distinguished subset
$\mathcal{A} \subseteq \mathcal{P}$ of \emph{admissible configurations}
satisfying the current constraint set $\mathcal{C}$. The \emph{actuality set}
$\mathcal{E} \subseteq \mathcal{P}$ is the set of configurations that have
been committed: extracted from $\mathcal{A}$ and rendered actual.
\end{definition}

\section{The Four Operators}

\begin{definition}{Pop}
$\mathrm{Pop}: \mathcal{A} \times \mathcal{P} \to \mathcal{P} \times \mathcal{E}$

\textbf{Signature:} Takes an admissible configuration $p \in \mathcal{A}$ and
the current possibility space $\mathcal{P}$; returns the contracted possibility
space $\mathcal{P}' = \mathcal{P} \setminus \{p\}$ and the updated actuality set
$\mathcal{E}' = \mathcal{E} \cup \{p\}$.

\textbf{Effect:} Extracts a committed event from possibility into actuality. All
alternatives to $p$ that are incompatible with $p$ given the active Bindings are
simultaneously removed from $\mathcal{A}$ (but not necessarily from $\mathcal{P}$,
since they may become relevant again if the Binding is released---though Bindings
themselves are irreversible).

\textbf{Irreversibility:} Once $\mathrm{Pop}(p)$ has been applied, $p$ cannot
be returned to $\mathcal{A}$. The event is committed.

\textbf{Example:} A Neanderthal builder selects a stalagmite fragment of length
$L = 0.47\,\text{m}$ from the set of available fragments. The fragment is
extracted from the possibility space of candidate elements and placed in the ring.
All fragments that would produce identical pitch are still available (they remain
in $\mathcal{P}$), but this specific fragment is now committed.
\end{definition}

\begin{definition}{Refuse}
$\mathrm{Refuse}: 2^\mathcal{P} \times \mathcal{P} \to \mathcal{P}$

\textbf{Signature:} Takes a set of inadmissible configurations $R \subseteq
\mathcal{P}$ and the current possibility space; returns $\mathcal{P}' =
\mathcal{P} \setminus R$.

\textbf{Effect:} Eliminates inadmissible trajectories from the possibility space
without committing to any alternative. Refuse is the admissibility filtering
operation that produces $\mathcal{A}(\mathcal{C})$ from the unconstrained space.

\textbf{Irreversibility:} Once refused, configurations cannot be re-admitted
within the same system-instance. The constraint that grounds the refusal may be
relaxed in a different context, but the Refuse operation itself leaves structural
residue.

\textbf{Example:} A squirrel, upon biting an object and finding it inedible,
removes it from its cache-candidate possibility space. The object remains in
the physical environment (it remains in the world-state) but is excluded from
the system's admissible foraging trajectories.
\end{definition}

\begin{definition}{Bind}
$\mathrm{Bind}: \mathcal{E} \times \mathcal{P} \to \mathcal{B}$

\textbf{Signature:} Takes a committed event $e \in \mathcal{E}$ and a prospective
event $p \in \mathcal{P}$; returns a dependency record $(e \to p) \in \mathcal{B}$
asserting that the occurrence of $e$ imposes constraints on the admissibility of $p$.

\textbf{Effect:} Establishes structural dependencies between events, constraining
future possibilities. Bind creates the dependency graph that constitutes the
provenance structure of the system.

\textbf{Irreversibility:} Active Bindings cannot be released within the same
system-instance. They accumulate the structural residue that makes the system's
history recoverable.

\textbf{Example:} In the Yarncrawler climate simulation, the Seattle tax-redirection
strategy is Bound to the Portland cooling-center expansion: the success of the
former constrains the admissible implementations of the latter. Prior decisions
structure the possibility space of future ones.
\end{definition}

\begin{definition}{Collapse}
$\mathrm{Collapse}: \mathcal{P} \to \mathcal{E}$

\textbf{Signature:} Takes the entire current possibility space $\mathcal{P}$
and commits it to a definite configuration, rendering all alternatives inaccessible.

\textbf{Effect:} Terminates the local trajectory of hypothesis refinement. Where
Pop extracts a single event, Collapse resolves an entire horizon of possibilities
into a committed configuration. After Collapse, no further Pop or Refuse operations
are possible on the committed domain.

\textbf{Irreversibility:} Collapse is the strongest irreversibility operation.
It commits everything. The structural residue of a Collapse is the entire history
of the possibility space that was resolved.

\textbf{Example:} At the end of the Bruniquel construction, the arrangement of
stalagmite fragments is Collapsed into a definite spatial configuration. The
alternative arrangements that were considered during construction are no longer
accessible to the same construction process---though they may become relevant to
a future reconstruction attempt.
\end{definition}

\section{Composition Rules}

Spherepop operators compose to produce \KES trajectories:

\[
  \Omega_{t+1} = \mathrm{Pop}^{k} \circ \mathrm{Refuse}^{m} \circ
  \mathrm{Bind}^{n} (\Omega_t)
\]

for appropriate choices of the exponent parameters. The ordering constraint is
that Bind must precede Pop (dependencies must be established before events are
committed under them), and Refuse must precede Pop (inadmissible options must
be eliminated before remaining options are selected). Collapse may only follow
when $|\mathcal{A}| = 1$ (the admissible space has contracted to a single option)
or as an executive operation terminating further refinement.

\section{The Bone Matrix Correspondence}

The seven Yarncrawler swarm-care axioms introduced in
Chapter~\ref{ch:yarncrawler_formal} correspond to Spherepop operation sequences:

\begin{center}
\begin{tabular}{ll}
\toprule
Swarm-care Axiom & Spherepop Sequence \\
\midrule
Local Action Only & $\mathrm{Refuse}(\text{non-local trajectories})$ \\
Trail Documentation & $\mathrm{Bind}(\text{action}, \text{log entry})$ \\
Autocatalytic Fix & $\mathrm{Bind}(\text{fix}, \text{future repair path})$ \\
Trail Reinforcement/Decay & $\mathrm{Pop}(\text{high-weight trail})$ \\
Node Hygiene & $\mathrm{Refuse}(\text{entropy-exceeding states})$ \\
Failure Isolation & $\mathrm{Collapse}(\text{quarantine zone})$ \\
Swarm Feedback & $\mathrm{Bind}(\text{precision-weighted update}, \text{global})$ \\
\bottomrule
\end{tabular}
\end{center}

This correspondence shows that the Yarncrawler axioms are not independent
engineering heuristics but derived consequences of operating Spherepop sequences
under the constraint that Markov blanket factorization is preserved throughout.
